\begin{explanationbox}
\textbf{\textcolor{CExplanation}{一句话直觉}}

只要下标和相同，项的和就相同。像对称的两组“脚标”不管具体是哪两项，只要下标总和一样，加起来的结果就固定。

\medskip
\textbf{\textcolor{CExplanation}{核心拆解}}
\begin{description}[style=nextline, leftmargin=2.8em, labelsep=0.5em, itemsep=0.45em]
    \item[\textbf{要点一（通项线性）：}] 等差数列通项 $a_n = a_1 + (n-1)d$ 是下标的线性函数。
    \item[\textbf{要点二（和消去 d）：}] 求和时 $a_m+a_n = 2a_1 + (m+n-2)d$，结果只依赖下标和 $m+n$。
    \item[\textbf{要点三（下标捆绑）：}] 只要 $m+n$ 固定，无论 $m,n$ 如何变化，和都不变。
\end{description}

\medskip
\textbf{\textcolor{CExplanation}{几何本质（直线中点）}}

将项 $(n, a_n)$ 看作直线上坐标为 $n$、纵坐标为 $a_n$ 的点。等差数列对应直线上等间距点。$m+n$ 相等意味着两对点的横坐标平均值相同，而和 $a_m+a_n$ 是纵坐标和，只依赖于下标和，因此和不变。

\medskip
\textbf{\textcolor{CExplanation}{代数意义（和消元）}}

结论将四项和的关系转化为下标和关系，避免了逐项计算。可以快速由部分和求其他项和，也可建立方程。

\medskip
\textbf{\textcolor{CExplanation}{考点价值（出题方向）}}
\begin{itemize}[leftmargin=*, itemsep=0.5em, topsep=0.4em]
    \item \textbf{考法一（快速求和）：} 已知某些项和，直接利用下标和相等口算另一组项和。
    \item \textbf{考法二（解方程）：} 设置下标关系，利用和相等等式建立方程求指定项或公差。
\end{itemize}

\medskip
\textbf{\textcolor{CExplanation}{顿悟点}}

\textbf{不是记住结论，而是记住“下标和”这个开关。} 看到下标和相等，立刻反应项和相等；看到下标和满足 $m+n=2k$，立刻想到等差中项。

\medskip
\textbf{\textcolor{CExplanation}{使用场景}}
\begin{itemize}[leftmargin=*, itemsep=0.5em, topsep=0.4em]
    \item \textbf{场景一（已知两项和）：} 例如已知 $a_2+a_5$，求 $a_3+a_4$。
    \item \textbf{场景二（下标对称）：} 题目出现“前几项和”、“中间项”时，通过构造下标和相等，快速关联。
\end{itemize}
\end{explanationbox}
